\section{\label{sec:intro}Introduction}

Variational quantum algorithms are among the main candidates for testing optimization routines on near-term quantum devices~\cite{Cerezo2021,Moll2018}.
The quantum approximate optimization algorithm (QAOA) prepares a parameterized state by alternating cost-phase and mixing unitaries, with the parameters optimized classically~\cite{Farhi2014,Blekos2024}. In its original form QAOA is suited to unconstrained binary optimization. Constrained problems are often treated by adding penalty terms to the cost Hamiltonian, which embeds the feasible problem into a larger Hilbert space and may assign probability to infeasible bit strings, with a performance that depends on the relative scale of objective and penalty~\cite{Hadfield2019,Blekos2024}.

Portfolio selection provides a constrained binary optimization problem of practical interest. Each binary variable records whether an asset is held, and a cardinality constraint fixes the number of selected assets. Quantum approaches have been studied using QAOA and related variational methods~\cite{Hodson2019,Slate2021}.

\sam{Recent work on multiclass portfolio optimization also uses parameterized Dicke-state ansatz to enforce class-wise fixed-cardinality constraints directly in the variational state, avoiding penalty calibration~\cite{Scursulim2026}.}

For fixed cardinality the ansatz can be built directly in the Hamming-weight sector of the allowed portfolios, which
avoids sampling infeasible configurations by construction. The quantum alternating operator ansatz permits mixers adapted to the feasible set~\cite{Hadfield2019}, and the quantum walk optimization algorithm (QWOA) generates the mixer from a continuous-time quantum walk on a graph whose vertices are the feasible configurations~\cite{Marsh2019,Marsh2020}. Different
graphs give different mixers: the feasible subspace is fixed, but the connectivity used to explore it is a free choice.

Our first contribution is to make that correspondence explicit. \sam{It is worth considering whether this contribution should be prioritized.}
Constraint-preserving mixers are usually written as Pauli sums and quantum walks as adjacency operators, and the identification between them is made case by case~\cite{Marsh2019, Marsh2020, Slate2021, Bartschi2020, Fuchs2022, Matwiejew2024}. Lemma~\ref{lem:dictionary} states it as a mixer--walk correspondence with
explicit hypotheses, so that it applies to any feasibility-preserving mixer,
and Remark~\ref{rem:failure} records what the mixer generates when the
hypotheses fail. The
Dicke-state plus Grover-mixer ansatz is exactly QWOA on the complete feasible
graph, so the Grover mixer of Ref.~\cite{Bartschi2020} and the complete-graph
mixer of Ref.~\cite{Slate2021} are the same operator; the equivalence is noted
in Ref.~\cite{Bridi2024} without the rescaling, which
Proposition~\ref{prop:grover-complete-qwoa} supplies as $\theta=M\beta/2$. And
the graph on which the ring
$XY$ mixer walks has vertex degree $2b(x)$, twice the number of maximal cyclic
blocks of ones in $x$, and is therefore not regular.

The second contribution is numerical. If every feasible mixer is a walk on a
graph, the graph is a design parameter, and the Johnson and complete graphs
are only two points in a large space. We therefore sweep the family
interpolating between them through unions of Johnson-scheme distance classes,
every member of which is regular and vertex-transitive, so that the uniform
feasible state remains an eigenvector of the mixer throughout. Performance is
not a trend in edge density: almost the whole family is statistically
indistinguishable from $J(n,k)$ across two orders of magnitude in degree, a
few unions sit apart, and the complete graph is an isolated outlier whose gap
widens with circuit depth.
That the complete graph is the worst member is expected:
Ref.~\cite{Bridi2024} proves that the Grover mixer is invariant under
permutations of the feasible states and hence blind to the structure of the
cost landscape, which by Proposition~\ref{prop:grover-complete-qwoa} applies to
the complete-graph mixer as well. What the sweep simulation adds is the interior of the
family, and in particular how early the degradation sets in.

Section~\ref{sec:preliminaries} fixes the portfolio problem, the
feasible-subspace ansatz, and the graph families used throughout.
    Section~\ref{sec:mixers-as-walks} gives the mixer--walk correspondence and applies it to the
    standard mixers. Section~\ref{sec:results} describes the data, the optimization protocol and reports the
    benchmarks and Sec.~\ref{sec:conclusion} concludes.